% ============================================================================ Anexos
% fuente: f9-remedicion-v2.md, f9-remedicion-completa-v2.md (auditorías, dev), f11-test-resultado.md §3,
% experiments/test-modelos-v2/{curve,operating}.csv (tablas generadas), configs/data/features_v1.yaml,
% configs/data/panel_seq_trips_v2_estaticas.yaml.
\section{Auditorías de fuga de información}
\label{anx:auditorias}

Además de las defensas de construcción (gap de 500\,km sobre la referencia corrida, partición por vehículo con los
duplicados unificados, variables solo hacia atrás y corte en el fin de extracción común), todo candidato se somete a tres
auditorías~\cite{kaufman2012}. Sea $s_{i,c}$ el puntaje fuera de pliegue y $\mathrm{AP}(\cdot)$ el PR-AUC por fila.
\begin{itemize}
  \item \textbf{(a0) Nulo global.} Se permutan las variables entre todas las filas, se reentrena y se exige
    $\mathrm{AP} \approx \bar y = 0{,}058$. Si no cae a la tasa base, hay información que no proviene de las variables.
  \item \textbf{\aprime{} Contribución del instante.} Sin reentrenar, se reemplaza cada puntaje por el promedio de su
    vehículo, $\bar s_i = \frac{1}{n_i}\sum_c s_{i,c}$, y se reporta $\mathrm{AP}(s) - \mathrm{AP}(\bar s)$: lo que el modelo
    sabía sobre \emph{cuándo}, más allá de \emph{qué vehículo}.
  \item \textbf{(b) Atajo de calendario.} Se agregan como variables las columnas de calendario (fechas de producción y de
    venta, temperatura ambiente) y se marca un atajo si el ROC sube más de $0{,}02$.
\end{itemize}

\begin{table}[!htbp]
\centering
\caption{Auditorías en desarrollo \tagdev{} (semilla 42 de la GRU; tasa base $0{,}058$).}
\label{tab:auditorias}
\begin{tabular}{@{}lrrrl@{}}
\toprule
Modelo & (a0) PR-AUC & \aprime{} & (b) $\Delta$ROC & Resultado \\
\midrule
GRU + viajes + estática (finalista) & 0,061 & $+0{,}001$ & $+0{,}004$ & aprueba \\
Survival stacking + motor y serie & 0,056 & $+0{,}025$ & $+0{,}009$ & aprueba \\
\bottomrule
\end{tabular}
\end{table}

Ninguno de los dos presenta fuga: con las variables permutadas, el PR-AUC cae a la tasa base. La contribución \aprime{} es
positiva en ambos, y el puntaje del finalista crece a medida que se acerca el evento (Figura~\ref{fig:score}). Ninguno
muestra el atajo de calendario.

\section{Resultados de todos los modelos}
\label{anx:modelos}

Cada modelo se entrenó con todo el conjunto de desarrollo (semilla 42, salvo la GRU, que es el ensamble de tres semillas) y
se evaluó con la misma métrica en los 103 vehículos de prueba. Estas evaluaciones se hicieron después de la evaluación del
finalista y son descriptivas: no se usaron para elegir.

\begin{table}[!htbp]
\centering
\caption{Detección $D(\alpha)$ (\%) en el conjunto de prueba según el presupuesto de falsas alarmas \tagtest{}. ``Prim.''
es $\bar D$, ec.~\eqref{eq:primario}.}
\label{tab:modelos-test}
\resizebox{\textwidth}{!}{% Generado por scripts/make_report_figures.py desde experiments/test-modelos-v2/curve.csv (no editar a mano).
\begin{tabular}{@{}lrrrrrrrrrl@{}}
\toprule
Modelo & 2\,\% & 5\,\% & 10\,\% & 15\,\% & 20\,\% & 25\,\% & 30\,\% & Prim. & AUC auto & dev (5 · 10 · 20\,\%) \\
\midrule
GRU ×3 (finalista) & 15,6 & 31,2 & 46,9 & 62,5 & 62,5 & 65,6 & 68,8 & \textbf{50,8} & 0,794 & 27 · 42 · 61 \\
CNN-LSTM ×3 & 3,1 & 15,6 & 56,2 & 56,2 & 68,8 & 68,8 & 68,8 & \textbf{49,2} & 0,790 & 22 · 41 · 62 \\
Survival stacking & 3,1 & 18,8 & 21,9 & 21,9 & 43,8 & 53,1 & 53,1 & \textbf{26,6} & 0,710 & 14 · 24 · 42 \\
LightGBM (53 features) & 0,0 & 15,6 & 25,0 & 28,1 & 34,4 & 53,1 & 56,2 & \textbf{25,8} & 0,725 & 18 · 28 · 42 \\
Logística L2 (53 features) & 0,0 & 6,2 & 12,5 & 31,2 & 34,4 & 46,9 & 53,1 & \textbf{21,1} & 0,679 & 15 · 23 · 37 \\
SS + efecto aleatorio (GPBoost) & 3,1 & 9,4 & 9,4 & 9,4 & 12,5 & 12,5 & 12,5 & \textbf{10,2} & 0,363 & 6 · 9 · 17 \\
\midrule
Azar, media / p95 (nulo de la GRU) & 3,4 / 9,4 & 7,0 / 15,6 & 13,7 / 25,0 & 18,6 / 31,2 & 24,4 / 40,6 & 28,6 / 43,8 & 33,9 / 50,0 & — & — & — \\
\bottomrule
\end{tabular}
}
\end{table}

\begin{table}[!htbp]
\centering
\caption{Con el umbral fijado en desarrollo \tagtest{}: detección y tasa de falsas alarmas realizada en prueba.}
\label{tab:operacion}
\begin{tabular}{@{}lll@{}}
\toprule
Modelo & Detección (\%) al 5 · 10 · 20\,\% & Falsas alarmas reales (\%) \\
\midrule
GRU ×3 (finalista) & 20 · 34 · 57 & 5,8 · 8,9 · 18,2 \\
Survival stacking & 13 · 20 · 33 & 4,6 · 8,3 · 15,6 \\
LightGBM (53 features) & 16 · 25 · 36 & 7,5 · 10,7 · 16,9 \\
Logística L2 (53 features) & 12 · 16 · 33 & 8,8 · 11,8 · 18,3 \\
SS + efecto aleatorio (GPBoost) & 8 · 11 · 14 & 6,5 · 11,3 · 21,5 \\
\bottomrule
\end{tabular}

\end{table}

\textbf{Comparación pareada en prueba} (bootstrap por vehículo, diferencia de $\bar D$ en puntos, IC 95\,\%): GRU $-$
survival stacking, $+24{,}2$ $[0{,}8;\,43{,}0]$.

\begin{table}[!htbp]
\centering
\caption{Modelos principales en desarrollo \tagdev{}: $D(\alpha)$ (\%), promedio de 3 repeticiones de la validación cruzada.}
\label{tab:dev}
\begin{tabular}{@{}lrrrr@{}}
\toprule
Modelo & AUC por vehículo & $\alpha=5\,\%$ & $10\,\%$ & $20\,\%$ \\
\midrule
GRU + viajes + estática, ×3 (finalista) & 0,82 & 27 & 42 & 61 \\
Survival stacking & 0,72 & 14 & 24 & 42 \\
LightGBM (53 variables) & 0,72 & 18 & 28 & 42 \\
Logística L2 & 0,71 & 15 & 23 & 37 \\
Nulo (historial conservado) & -- & 5 & 11 & 21 \\
\bottomrule
\end{tabular}
\end{table}

\section{Diccionario de variables}
\label{anx:features}

Las 53 variables agregadas se calculan sobre la ventana $(c - W, c]$ según la ec.~\eqref{eq:agregado}, como
\texttt{agregador(variable derivada)}. Las variables con prefijo \texttt{aux\_} se conservan para auditoría y no entran en
ningún modelo. Agregadores: \texttt{per\_1000km}, conteo por 1.000\,km; \texttt{per\_day}, por día; \texttt{half\_*},
segunda mitad de la ventana menos la primera; \texttt{gap\_*}, distancia entre eventos; \texttt{any}, ocurrencia.

{\scriptsize
% Generado por scripts/make_report_figures.py desde configs/data/features_v1.yaml (no editar a mano).
% 53 features de modelo + 3 aux.
\begin{longtable}{@{}>{\raggedright\arraybackslash}p{5.4cm}l>{\raggedright\arraybackslash}p{4.6cm}>{\raggedright\arraybackslash}p{3.2cm}@{}}
\toprule
Feature & Fuente & Columna derivada & Agregador \\
\endhead
\midrule
\multicolumn{4}{@{}l}{\textbf{A · Térmica y trayectos cortos}} \\
\texttt{feat\_\allowbreak{}idle\_\allowbreak{}frac} & trips & \texttt{idle} & \texttt{mean} \\
\texttt{feat\_\allowbreak{}idle\_\allowbreak{}per\_\allowbreak{}1000km} & trips & \texttt{idle} & \texttt{per\_\allowbreak{}1000km} \\
\texttt{feat\_\allowbreak{}idle\_\allowbreak{}frac\_\allowbreak{}trend} & trips & \texttt{idle} & \texttt{half\_\allowbreak{}mean\_\allowbreak{}delta} \\
\texttt{feat\_\allowbreak{}short\_\allowbreak{}trip\_\allowbreak{}frac\_\allowbreak{}5km} & trips & \texttt{short\_\allowbreak{}trip\_\allowbreak{}moving} & \texttt{mean} \\
\texttt{feat\_\allowbreak{}trips\_\allowbreak{}below\_\allowbreak{}regime\_\allowbreak{}temp\_\allowbreak{}frac} & trips & \texttt{below\_\allowbreak{}regime} & \texttt{mean} \\
\texttt{feat\_\allowbreak{}below\_\allowbreak{}regime\_\allowbreak{}moving\_\allowbreak{}frac} & trips & \texttt{below\_\allowbreak{}regime\_\allowbreak{}moving} & \texttt{mean} \\
\texttt{feat\_\allowbreak{}below\_\allowbreak{}regime\_\allowbreak{}frac\_\allowbreak{}trend} & trips & \texttt{below\_\allowbreak{}regime} & \texttt{half\_\allowbreak{}mean\_\allowbreak{}delta} \\
\texttt{feat\_\allowbreak{}engine\_\allowbreak{}temp\_\allowbreak{}avg\_\allowbreak{}median} & trips & \texttt{EngineTemperatureAvg\_\allowbreak{}moving} & \texttt{median} \\
\texttt{feat\_\allowbreak{}engine\_\allowbreak{}temp\_\allowbreak{}max\_\allowbreak{}median} & trips & \texttt{EngineTemperatureMax\_\allowbreak{}moving} & \texttt{median} \\
\texttt{feat\_\allowbreak{}engine\_\allowbreak{}temp\_\allowbreak{}amplitude\_\allowbreak{}mean} & trips & \texttt{engine\_\allowbreak{}temp\_\allowbreak{}amplitude\_\allowbreak{}moving} & \texttt{mean} \\
\texttt{feat\_\allowbreak{}coolant\_\allowbreak{}temp\_\allowbreak{}end\_\allowbreak{}mean} & trips & \texttt{CoolantTemperatureEnd\_\allowbreak{}moving} & \texttt{mean} \\
\texttt{feat\_\allowbreak{}coolant\_\allowbreak{}temp\_\allowbreak{}start\_\allowbreak{}median} & trips & \texttt{CoolantTemperatureStart} & \texttt{median} \\
\texttt{feat\_\allowbreak{}cold\_\allowbreak{}start\_\allowbreak{}frac} & trips & \texttt{cold\_\allowbreak{}start} & \texttt{mean} \\
\texttt{feat\_\allowbreak{}chained\_\allowbreak{}trip\_\allowbreak{}frac} & trips & \texttt{chained} & \texttt{mean} \\
\texttt{feat\_\allowbreak{}trip\_\allowbreak{}distance\_\allowbreak{}median\_\allowbreak{}km} & trips & \texttt{trip\_\allowbreak{}km\_\allowbreak{}moving} & \texttt{median} \\
\texttt{feat\_\allowbreak{}trip\_\allowbreak{}distance\_\allowbreak{}p25\_\allowbreak{}km} & trips & \texttt{trip\_\allowbreak{}km\_\allowbreak{}moving} & \texttt{p25} \\
\texttt{feat\_\allowbreak{}trip\_\allowbreak{}duration\_\allowbreak{}median\_\allowbreak{}min} & trips & \texttt{trip\_\allowbreak{}duration\_\allowbreak{}min\_\allowbreak{}moving} & \texttt{median} \\
\midrule
\multicolumn{4}{@{}l}{\textbf{B · Regeneración (DPF)}} \\
\texttt{feat\_\allowbreak{}regenerations\_\allowbreak{}per\_\allowbreak{}1000km} & trips & \texttt{regen\_\allowbreak{}drop} & \texttt{per\_\allowbreak{}1000km} \\
\texttt{feat\_\allowbreak{}regenerations\_\allowbreak{}trend\_\allowbreak{}per\_\allowbreak{}1000km} & trips & \texttt{regen\_\allowbreak{}drop} & \texttt{half\_\allowbreak{}count\_\allowbreak{}delta\_\allowbreak{}per\_\allowbreak{}1000km} \\
\texttt{feat\_\allowbreak{}distance\_\allowbreak{}between\_\allowbreak{}regen\_\allowbreak{}mean\_\allowbreak{}km} & trips & \texttt{regen\_\allowbreak{}drop} & \texttt{gap\_\allowbreak{}mean\_\allowbreak{}km} \\
\texttt{feat\_\allowbreak{}distance\_\allowbreak{}between\_\allowbreak{}regen\_\allowbreak{}trend} & trips & \texttt{regen\_\allowbreak{}drop} & \texttt{gap\_\allowbreak{}trend\_\allowbreak{}km} \\
\texttt{feat\_\allowbreak{}km\_\allowbreak{}since\_\allowbreak{}last\_\allowbreak{}regen} & trips & \texttt{regen\_\allowbreak{}drop} & \texttt{km\_\allowbreak{}since\_\allowbreak{}last} \\
\texttt{feat\_\allowbreak{}regen\_\allowbreak{}residual\_\allowbreak{}mean} & trips & \texttt{regen\_\allowbreak{}residual} & \texttt{mean} \\
\texttt{feat\_\allowbreak{}regen\_\allowbreak{}start\_\allowbreak{}level\_\allowbreak{}mean} & trips & \texttt{regen\_\allowbreak{}start\_\allowbreak{}level} & \texttt{mean} \\
\texttt{feat\_\allowbreak{}dpf\_\allowbreak{}end\_\allowbreak{}mean} & trips & \texttt{AirRegenerationEnd} & \texttt{mean} \\
\texttt{feat\_\allowbreak{}dpf\_\allowbreak{}end\_\allowbreak{}max} & trips & \texttt{AirRegenerationEnd} & \texttt{max} \\
\texttt{feat\_\allowbreak{}dpf\_\allowbreak{}end\_\allowbreak{}slope\_\allowbreak{}per\_\allowbreak{}1000km} & trips & \texttt{AirRegenerationEnd} & \texttt{half\_\allowbreak{}slope\_\allowbreak{}per\_\allowbreak{}1000km} \\
\texttt{feat\_\allowbreak{}dpf\_\allowbreak{}positive\_\allowbreak{}delta\_\allowbreak{}frac} & trips & \texttt{dpf\_\allowbreak{}positive\_\allowbreak{}delta\_\allowbreak{}moving} & \texttt{mean} \\
\texttt{feat\_\allowbreak{}dpf\_\allowbreak{}saturated\_\allowbreak{}frac} & trips & \texttt{dpf\_\allowbreak{}saturated\_\allowbreak{}end} & \texttt{mean} \\
\texttt{feat\_\allowbreak{}filter\_\allowbreak{}abnormal\_\allowbreak{}end\_\allowbreak{}frac} & trips & \texttt{filter\_\allowbreak{}abnormal\_\allowbreak{}end} & \texttt{mean} \\
\texttt{feat\_\allowbreak{}filter\_\allowbreak{}cleaning\_\allowbreak{}auto\_\allowbreak{}end\_\allowbreak{}frac} & trips & \texttt{filter\_\allowbreak{}cleaning\_\allowbreak{}auto\_\allowbreak{}end} & \texttt{mean} \\
\texttt{feat\_\allowbreak{}manual\_\allowbreak{}regen\_\allowbreak{}ever} & trips & \texttt{filter\_\allowbreak{}manual} & \texttt{any} \\
\midrule
\multicolumn{4}{@{}l}{\textbf{C · Uso}} \\
\texttt{feat\_\allowbreak{}speed\_\allowbreak{}kmh\_\allowbreak{}mean} & trips & \texttt{speed\_\allowbreak{}kmh\_\allowbreak{}moving} & \texttt{mean} \\
\texttt{feat\_\allowbreak{}trips\_\allowbreak{}below\_\allowbreak{}30kmh\_\allowbreak{}frac} & trips & \texttt{urban\_\allowbreak{}moving} & \texttt{mean} \\
\texttt{feat\_\allowbreak{}moving\_\allowbreak{}trips\_\allowbreak{}per\_\allowbreak{}1000km} & trips & \texttt{moving} & \texttt{per\_\allowbreak{}1000km} \\
\texttt{feat\_\allowbreak{}km\_\allowbreak{}per\_\allowbreak{}day} & trips & \texttt{trip\_\allowbreak{}km} & \texttt{per\_\allowbreak{}day} \\
\texttt{feat\_\allowbreak{}trips\_\allowbreak{}per\_\allowbreak{}day} & trips & \texttt{moving} & \texttt{per\_\allowbreak{}day} \\
\texttt{feat\_\allowbreak{}hours\_\allowbreak{}between\_\allowbreak{}trips\_\allowbreak{}median} & trips & \texttt{soak\_\allowbreak{}min} & \texttt{median} \\
\texttt{aux\_\allowbreak{}air\_\allowbreak{}temp\_\allowbreak{}avg} & trips & \texttt{AirTemperatureAvg} & \texttt{mean} \\
\texttt{aux\_\allowbreak{}air\_\allowbreak{}temp\_\allowbreak{}min} & trips & \texttt{AirTemperatureMin} & \texttt{mean} \\
\midrule
\multicolumn{4}{@{}l}{\textbf{D · Severidad (mensajes, aceite, consumo)}} \\
\texttt{feat\_\allowbreak{}msg\_\allowbreak{}full\_\allowbreak{}per\_\allowbreak{}1000km} & signals & \texttt{msg\_\allowbreak{}full} & \texttt{per\_\allowbreak{}1000km} \\
\texttt{feat\_\allowbreak{}msg\_\allowbreak{}overloaded\_\allowbreak{}per\_\allowbreak{}1000km} & signals & \texttt{msg\_\allowbreak{}overloaded} & \texttt{per\_\allowbreak{}1000km} \\
\texttt{feat\_\allowbreak{}msg\_\allowbreak{}overloaded\_\allowbreak{}ever} & signals & \texttt{msg\_\allowbreak{}overloaded} & \texttt{any} \\
\texttt{feat\_\allowbreak{}msg\_\allowbreak{}over\_\allowbreak{}limit\_\allowbreak{}per\_\allowbreak{}1000km} & signals & \texttt{msg\_\allowbreak{}over\_\allowbreak{}limit} & \texttt{per\_\allowbreak{}1000km} \\
\texttt{feat\_\allowbreak{}msg\_\allowbreak{}over\_\allowbreak{}limit\_\allowbreak{}ever} & signals & \texttt{msg\_\allowbreak{}over\_\allowbreak{}limit} & \texttt{any} \\
\texttt{feat\_\allowbreak{}msg\_\allowbreak{}at\_\allowbreak{}limit\_\allowbreak{}ever} & signals & \texttt{msg\_\allowbreak{}at\_\allowbreak{}limit} & \texttt{any} \\
\texttt{feat\_\allowbreak{}msg\_\allowbreak{}cleaning\_\allowbreak{}auto\_\allowbreak{}per\_\allowbreak{}1000km} & signals & \texttt{msg\_\allowbreak{}cleaning\_\allowbreak{}auto} & \texttt{per\_\allowbreak{}1000km} \\
\texttt{feat\_\allowbreak{}msg\_\allowbreak{}cleaning\_\allowbreak{}auto\_\allowbreak{}trend\_\allowbreak{}per\_\allowbreak{}1000km} & signals & \texttt{msg\_\allowbreak{}cleaning\_\allowbreak{}auto} & \texttt{half\_\allowbreak{}count\_\allowbreak{}delta\_\allowbreak{}per\_\allowbreak{}1000km} \\
\texttt{feat\_\allowbreak{}msg\_\allowbreak{}abnormal\_\allowbreak{}frac} & signals & \texttt{msg\_\allowbreak{}abnormal} & \texttt{mean} \\
\texttt{feat\_\allowbreak{}msgs\_\allowbreak{}per\_\allowbreak{}1000km} & signals & \texttt{msg\_\allowbreak{}any} & \texttt{per\_\allowbreak{}1000km} \\
\texttt{feat\_\allowbreak{}oil\_\allowbreak{}life\_\allowbreak{}mean} & trips & \texttt{oil\_\allowbreak{}life} & \texttt{mean} \\
\texttt{feat\_\allowbreak{}oil\_\allowbreak{}life\_\allowbreak{}slope\_\allowbreak{}per\_\allowbreak{}1000km} & trips & \texttt{oil\_\allowbreak{}life} & \texttt{half\_\allowbreak{}slope\_\allowbreak{}per\_\allowbreak{}1000km} \\
\texttt{feat\_\allowbreak{}fuel\_\allowbreak{}pct\_\allowbreak{}per\_\allowbreak{}100km\_\allowbreak{}median} & trips & \texttt{fuel\_\allowbreak{}pct\_\allowbreak{}per\_\allowbreak{}100km} & \texttt{median} \\
\texttt{aux\_\allowbreak{}regen\_\allowbreak{}marker\_\allowbreak{}per\_\allowbreak{}1000km} & signals & \texttt{regen\_\allowbreak{}marker} & \texttt{per\_\allowbreak{}1000km} \\
\midrule
\multicolumn{4}{@{}l}{\textbf{Control de ventana (los emite \texttt{windows.py})}} \\
\texttt{feat\_n\_trips\_window} & trips & \texttt{*} & \texttt{count} \\
\texttt{feat\_window\_km\_covered} & trips & \texttt{OdometerTripEnd} & \texttt{rango} \\
\bottomrule
\end{longtable}

}

\textbf{Canales de la representación secuencial}~\eqref{eq:secuencia}. La ventana se divide en 20 tramos de 50\,km;
``zero'' rellena con cero un tramo sin datos y ``carry'' arrastra el último valor dentro de la ventana. Se agregan país, motor
y serie en codificación \emph{one-hot}.

\begin{center}
\footnotesize
% Generado por scripts/make_report_figures.py desde configs/data/panel_seq_trips_v2_estaticas.yaml (no editar a mano).
% ventana 1000 km en bins de 50 km = 20 tramos.
\begin{tabular}{@{}lllll@{}}
\toprule
Canal & Fuente & Columna & Agregado por tramo & Relleno \\
\midrule
\texttt{n\_records} & signals & \texttt{*} & count & zero · log1p \\
\texttt{observed} & signals & \texttt{*} & any & zero \\
\texttt{acum\_mean} & signals & \texttt{Acumulation} & mean & carry \\
\texttt{acum\_max} & signals & \texttt{Acumulation} & max & carry \\
\texttt{acum\_rise} & signals & \texttt{acum\_rise} & sum & zero \\
\texttt{regen\_drops} & signals & \texttt{acum\_regen\_drop} & sum & zero \\
\texttt{msg\_full\_frac} & signals & \texttt{msg\_full} & mean & zero \\
\texttt{msg\_severe\_frac} & signals & \texttt{msg\_severe} & mean & zero \\
\texttt{msg\_cleaning\_auto\_frac} & signals & \texttt{msg\_cleaning\_auto} & mean & zero \\
\texttt{n\_trips} & trips & \texttt{*} & count & zero · log1p \\
\texttt{idle\_frac} & trips & \texttt{idle} & mean & carry \\
\texttt{below\_regime\_frac} & trips & \texttt{below\_regime\_moving} & mean & carry \\
\texttt{speed\_kmh} & trips & \texttt{speed\_kmh\_moving} & mean & carry \\
\texttt{trip\_km} & trips & \texttt{trip\_km\_moving} & mean & carry \\
\texttt{trip\_duration\_min} & trips & \texttt{trip\_duration\_min\_moving} & mean & carry \\
\bottomrule
\end{tabular}

\end{center}
